\centerline{\bf \Huge Матанализ I}

\begin{flushright}
        \scriptsize{Советский ребёнок знал это ещё в утробе!}
\end{flushright}

\problem{1}
(а) \textbf{Лемма о вложенных отрезках.} Пусть $\left[a_1, b_1\right] \supset\left[a_2, b_2\right] \supset \ldots$ - бесконечная последовательность вложенных отрезков. Тогда существует число $c$, принадлежащее всем отрезкам, т.е. $\forall i \in \mathbb{N}$ выполнено $c \in\left[a_i, b_i\right]$.

(b) Докажите, что всякая ограниченная с двух сторон последовательность имеет сходящуюся подпоследовательность.


\problem{2}
\textbf{Арифметика пределов.}
Пусть $\left(a_n\right)_{n=1}^{\infty}$ и $\left(b_n\right)_{n=1}^{\infty}$ - две сходящиеся к $A$ и $B$ соответственно последовательности. Тогда

(a) последовательность $\left(a_n+b_n\right)_{n=1}^{\infty}$ сходится к $A+B$; другими словами

$$
\lim _{n \rightarrow \infty}\left(a_n+b_n\right)=\lim _{n \rightarrow \infty} a_n+\lim _{n \rightarrow \infty} b_n .
$$

(b) последовательность $\left(a_n b_n\right)_{n=1}^{\infty}$ сходится к $A B$; другими словами

$$
\lim _{n \rightarrow \infty}\left(a_n b_n\right)=\lim _{n \rightarrow \infty} a_n \cdot \lim _{n \rightarrow \infty} b_n .
$$

(c) пусть $B \neq 0$, а также $b_n \neq 0$ при $n \geqslant m$. Тогда последовательность $\left(a_n / b_n\right)_{n=m}^{\infty}$ сходится к $A / B$; другими словами

$$
\lim _{n \rightarrow \infty} \frac{a_n}{b_n}=\frac{\lim _{n \rightarrow \infty} a_n}{\lim _{n \rightarrow \infty} b_n} .
$$